\SourceLeaf{392}%
\[
\begin{aligned}
\mathrm{V.}\quad&2aa\sin2(A''-A')+2a'a'\sin2(A-A'')+2a''a''\sin2(A'-A)\\
&\quad+t[\sin2(A''-A')+\sin2(A-A'')+\sin2(A'-A)]\\
&\quad-4ar\cos(A-\varphi)\sin2(A''-A')\\
&\quad-4a'r\cos(A'-\varphi)\sin2(A-A'')\\
&\quad-4a''r\cos(A''-\varphi)\sin2(A'-A)=0
\end{aligned}
\]
\par
Die Summe, mit der \(t\) hier multipliziert wird, lässt sich auch durch
\[
4\sin(A''-A')\sin(A''-A)\sin(A'-A)
\]
ausdrücken.

Behandelt man die Gleichungen I, II, IV auf ähnliche Weise, erhält man eine entsprechende Gleichung VI, die sich von V nur dadurch unterscheidet, dass \(a'',A''\) durch \(a''',A'''\) ersetzt wird. Eliminiert man aus den beiden Gleichungen V und VI die Größe \(t\), so sieht man leicht, dass sich eine Gleichung der Form
\[
\mathrm{VII.}\qquad B+Cr\cos\varphi+Dr\sin\varphi=0
\]
ergibt, wobei \(B,C,D\) bekannte Größen bezeichnen. Ihre Werte lassen sich leicht bestimmen, doch werden wir bald zeigen, wie sich diese Rechnung vermeiden lässt. Aus Gleichung VII geht hervor, dass der Mittelpunkt jeder Ellipse, die die vier Seiten unseres Vierecks berührt, auf einer Geraden liegt. Diese ist gegen die Abszissenachse um einen Winkel mit der Tangente \(=-\frac{C}{D}\) geneigt, und der Schnittpunkt hat die Abszisse \(-\frac{B}{C}\). Die Lage dieser Geraden lässt sich aber viel leichter mit folgender Überlegung bestimmen. Eine Diagonale des Vierecks kann man als eine verschwindende Ellipse auffassen, die die Seiten des Vierecks berührt. Ihr Mittelpunkt liegt dann offenbar in der Mitte der Diagonale. Daraus folgt leicht, dass die oben genannte Gerade, der geometrische Ort der Mittelpunkte aller Ellipsen, die die vier Seiten des Vierecks berühren, die Verbindungslinie der Mittelpunkte beider Diagonalen sein muss. Sie lässt sich demnach leicht finden. Dazu füge ich noch zwei Bemerkungen hinzu.

1) Fallen beide Mittelpunkte zusammen (in diesem Fall ist das Viereck ein Parallelogramm), so entfällt freilich diese Bestimmung der
